% Lösungen Einheit 4 von 4 – Scharen am Körper (zusammenbau v0.1)
\einheitenkopf{Einheit 4 von 4 · Scharen am Körper}

\erg{16}{a) 0, 3, 6 und 9 – die Werte von $x_1 + x_2 + x_3$ an den Ecken \quad b) $E_2: x_1 + x_2 + x_3 = 2$; Dreieck mit den Ecken $(2 | 0 | 0)$, $(0 | 2 | 0)$, $(0 | 0 | 2)$ auf den Kanten durch den Ursprung \quad c) $R(2 | 4 | 4)$: nur die Ecke $(4 | 4 | 4)$ liefert mehr als $10$; auf ihrer Kante zu $(0 | 4 | 4)$ gilt $x_1 + 8 = 10$ \quad d) Übergangswerte an den Ecken: $k = 2; 4$; drei Ecken für $0 < k \le 2$, sechs Ecken für $2 < k < 4$, drei Ecken für $4 \le k < 6$ \quad e) Übergänge $k = 0; 2; 6; 8$; größter Bereich $]2; 6[$ – dort schneidet $E_k$ die vier Kanten in $x_1$-Richtung \quad f) $x_3 = 0$: $k x_1 + x_2 = 4$, also $g_k: \vec{x} = (0 | 4 | 0) + r \cdot (1 | -k | 0)$ \quad g) größter Wert von $x_1 + 3x_2 + x_3$ über den Ecken: $4 + 6 + 6 = 16$ in $(4 | 2 | 6)$; für $k > 16$ liegt der ganze Quader auf einer Seite von $E_k$, also $k_0 = 16$ (nicht die größte Koordinate $6$)}
\erg{17}{a) Fehler: die größte Koordinate statt des größten Werts von $x_1 + 2x_2 + x_3$ genommen. Richtig: in $(3 | 2 | 4)$ ist der Wert $3 + 4 + 4 = 11$, also $k_0 = 11$ \quad b) Zwischen zwei Übergängen läuft die Ebene durch keine Ecke: jede Ecke bleibt auf ihrer Seite, also werden dieselben Kanten geschnitten und die Eckenzahl bleibt gleich.}
